% GENERATED FILE - do not hand-edit.
% Source of truth: data/calculus-reviewed/t-2.json in the zenith-website repo.
% Regenerate with: python3 scripts/build_calculus_tex.py
\documentclass[11pt]{article}
\usepackage[utf8]{inputenc}
\usepackage[T1]{fontenc}
\usepackage{amsmath,amssymb}
\usepackage{graphicx}
\usepackage[margin=1in]{geometry}
\graphicspath{{../..//}}
\setlength{\parindent}{0pt}
\setlength{\parskip}{0.6em}
\begin{document}
\section*{Chapter 2: Limits and Continuity --- Test T2}

\bigskip\noindent\textbf{T1.}\ $\displaystyle\lim_{x\to 3}\dfrac{x^2-9}{x^2+9}$ is
\begin{list}{}{\setlength{\leftmargin}{3em}\setlength{\labelwidth}{2.5em}\setlength{\itemsep}{0pt}\setlength{\topsep}{2pt}}
\item[(A)] $1$
\item[(B)] $0$
\item[(C)] $-\dfrac{1}{3}$
\item[(D)] $\infty$
\end{list}

\bigskip\noindent\textbf{T2.}\ $\displaystyle\lim_{x\to\infty}\dfrac{9-x^2}{x^2-4}$ is
\begin{list}{}{\setlength{\leftmargin}{3em}\setlength{\labelwidth}{2.5em}\setlength{\itemsep}{0pt}\setlength{\topsep}{2pt}}
\item[(A)] $1$
\item[(B)] $0$
\item[(C)] $-1$
\item[(D)] $\infty$
\end{list}

\bigskip\noindent\textbf{T3.}\ $\displaystyle\lim_{x\to 4}\dfrac{x-4}{x^2-3x-4}$ is
\begin{list}{}{\setlength{\leftmargin}{3em}\setlength{\labelwidth}{2.5em}\setlength{\itemsep}{0pt}\setlength{\topsep}{2pt}}
\item[(A)] $0$
\item[(B)] $1$
\item[(C)] $\dfrac{1}{5}$
\item[(D)] $\infty$
\end{list}

\bigskip\noindent\textbf{T4.}\ $\displaystyle\lim_{x\to 0}\dfrac{5x}{x}$ is
\begin{list}{}{\setlength{\leftmargin}{3em}\setlength{\labelwidth}{2.5em}\setlength{\itemsep}{0pt}\setlength{\topsep}{2pt}}
\item[(A)] $5$
\item[(B)] $0$
\item[(C)] $-5$
\item[(D)] nonexistent
\end{list}

\bigskip\noindent\textbf{T5.}\ $\displaystyle\lim_{x\to 4}\dfrac{x^3-64}{x^2-16}$ is
\begin{list}{}{\setlength{\leftmargin}{3em}\setlength{\labelwidth}{2.5em}\setlength{\itemsep}{0pt}\setlength{\topsep}{2pt}}
\item[(A)] $2$
\item[(B)] $4$
\item[(C)] $6$
\item[(D)] $\infty$
\end{list}

\bigskip\noindent\textbf{T6.}\ $\displaystyle\lim_{x\to\infty}\dfrac{9-x^2}{9x^2-2x-4}$ is
\begin{list}{}{\setlength{\leftmargin}{3em}\setlength{\labelwidth}{2.5em}\setlength{\itemsep}{0pt}\setlength{\topsep}{2pt}}
\item[(A)] $-3$
\item[(B)] $-\dfrac{1}{9}$
\item[(C)] $1$
\item[(D)] $3$
\end{list}

\bigskip\noindent\textbf{T7.}\ $\displaystyle\lim_{x\to-\infty}\dfrac{3x^3+40}{9x^2+6x+5}$ is
\begin{list}{}{\setlength{\leftmargin}{3em}\setlength{\labelwidth}{2.5em}\setlength{\itemsep}{0pt}\setlength{\topsep}{2pt}}
\item[(A)] $-\infty$
\item[(B)] $-1$
\item[(C)] $\dfrac{1}{3}$
\item[(D)] $\infty$
\end{list}

\bigskip\noindent\textbf{T8.}\ $\displaystyle\lim_{x\to\infty}\dfrac{5x^2+50}{x^3-125}$ is
\begin{list}{}{\setlength{\leftmargin}{3em}\setlength{\labelwidth}{2.5em}\setlength{\itemsep}{0pt}\setlength{\topsep}{2pt}}
\item[(A)] $0$
\item[(B)] $1$
\item[(C)] $5$
\item[(D)] $\infty$
\end{list}

\bigskip\noindent\textbf{T9.}\ $\displaystyle\lim_{x\to\infty}\dfrac{3^{-x}}{3^x}$ is
\begin{list}{}{\setlength{\leftmargin}{3em}\setlength{\labelwidth}{2.5em}\setlength{\itemsep}{0pt}\setlength{\topsep}{2pt}}
\item[(A)] $-1$
\item[(B)] $1$
\item[(C)] $0$
\item[(D)] $\infty$
\end{list}

\bigskip\noindent\textbf{T10.}\ $\displaystyle\lim_{x\to-\infty}\dfrac{3^{-x}}{3^x}$ is
\begin{list}{}{\setlength{\leftmargin}{3em}\setlength{\labelwidth}{2.5em}\setlength{\itemsep}{0pt}\setlength{\topsep}{2pt}}
\item[(A)] $-1$
\item[(B)] $1$
\item[(C)] $0$
\item[(D)] $\infty$
\end{list}

\bigskip\noindent\textbf{T11.}\ $\displaystyle\lim_{x\to 0}\dfrac{\sin 7x}{x}$
\begin{list}{}{\setlength{\leftmargin}{3em}\setlength{\labelwidth}{2.5em}\setlength{\itemsep}{0pt}\setlength{\topsep}{2pt}}
\item[(A)] $=\dfrac{1}{7}$
\item[(B)] $=1$
\item[(C)] $=7$
\item[(D)] does not exist
\end{list}

\bigskip\noindent\textbf{T12.}\ $\displaystyle\lim_{x\to 0}\dfrac{\sin 3x}{4x}$
\begin{list}{}{\setlength{\leftmargin}{3em}\setlength{\labelwidth}{2.5em}\setlength{\itemsep}{0pt}\setlength{\topsep}{2pt}}
\item[(A)] $=\dfrac{3}{4}$
\item[(B)] $=1$
\item[(C)] $=\dfrac{4}{3}$
\item[(D)] does not exist
\end{list}

\bigskip\noindent\textbf{T13.}\ The graph of $y = 2\arctan x$ has
\begin{list}{}{\setlength{\leftmargin}{3em}\setlength{\labelwidth}{2.5em}\setlength{\itemsep}{0pt}\setlength{\topsep}{2pt}}
\item[(A)] vertical asymptotes at $x = 0$ and $x = \pi$
\item[(B)] horizontal asymptotes at $y = \pm\pi$
\item[(C)] horizontal asymptotes at $y = 0$ and $y = \pi$
\item[(D)] vertical asymptotes at $x = \pm\dfrac{\pi}{2}$
\end{list}

\bigskip\noindent\textbf{T14.}\ The graph of $y = \dfrac{x^2-16}{4x-16}$ has
\begin{list}{}{\setlength{\leftmargin}{3em}\setlength{\labelwidth}{2.5em}\setlength{\itemsep}{0pt}\setlength{\topsep}{2pt}}
\item[(A)] a vertical asymptote at $x = 4$
\item[(B)] a horizontal asymptote at $y = \dfrac{1}{4}$
\item[(C)] a removable discontinuity at $x = 4$
\item[(D)] an infinite discontinuity at $x = 4$
\end{list}

\bigskip\noindent\textbf{T15.}\ $\displaystyle\lim_{x\to 0}\dfrac{\sin x}{x^2+5x}$ is
\begin{list}{}{\setlength{\leftmargin}{3em}\setlength{\labelwidth}{2.5em}\setlength{\itemsep}{0pt}\setlength{\topsep}{2pt}}
\item[(A)] $1$
\item[(B)] $\dfrac{1}{5}$
\item[(C)] $5$
\item[(D)] $\infty$
\end{list}

\bigskip\noindent\textbf{T16.}\ $\displaystyle\lim_{x\to 0}\cos\dfrac{1}{x}$ is
\begin{list}{}{\setlength{\leftmargin}{3em}\setlength{\labelwidth}{2.5em}\setlength{\itemsep}{0pt}\setlength{\topsep}{2pt}}
\item[(A)] $\infty$
\item[(B)] $1$
\item[(C)] $-1$
\item[(D)] nonexistent
\end{list}

\bigskip\noindent\textbf{T17.}\ Which statement is true about the curve $y = \dfrac{4x^2+16}{3+11x-4x^2}$?
\begin{list}{}{\setlength{\leftmargin}{3em}\setlength{\labelwidth}{2.5em}\setlength{\itemsep}{0pt}\setlength{\topsep}{2pt}}
\item[(A)] The line $x = -\dfrac{1}{4}$ is a vertical asymptote.
\item[(B)] The line $x = 1$ is a vertical asymptote.
\item[(C)] The line $y = -\dfrac{1}{4}$ is a horizontal asymptote.
\item[(D)] The line $y = 4$ is a horizontal asymptote.
\end{list}

\bigskip\noindent\textbf{T18.}\ $\displaystyle\lim_{x\to\infty}\dfrac{3x^2+1}{(3-x)(3+x)}$ is
\begin{list}{}{\setlength{\leftmargin}{3em}\setlength{\labelwidth}{2.5em}\setlength{\itemsep}{0pt}\setlength{\topsep}{2pt}}
\item[(A)] $-3$
\item[(B)] $1$
\item[(C)] $3$
\item[(D)] nonexistent
\end{list}

\bigskip\noindent\textbf{T19.}\ $\displaystyle\lim_{x\to 0}\dfrac{|3x|}{x}$ is
\begin{list}{}{\setlength{\leftmargin}{3em}\setlength{\labelwidth}{2.5em}\setlength{\itemsep}{0pt}\setlength{\topsep}{2pt}}
\item[(A)] $0$
\item[(B)] $3$
\item[(C)] $-3$
\item[(D)] nonexistent
\end{list}

\bigskip\noindent\textbf{T20.}\ $\displaystyle\lim_{x\to\infty} x\sin\dfrac{4}{x}$ is
\begin{list}{}{\setlength{\leftmargin}{3em}\setlength{\labelwidth}{2.5em}\setlength{\itemsep}{0pt}\setlength{\topsep}{2pt}}
\item[(A)] $0$
\item[(B)] $-4$
\item[(C)] $4$
\item[(D)] nonexistent
\end{list}

\bigskip\noindent\textbf{T21.}\ $\displaystyle\lim_{x\to 3\pi}\dfrac{\sin(3\pi-x)}{3\pi-x}$ is
\begin{list}{}{\setlength{\leftmargin}{3em}\setlength{\labelwidth}{2.5em}\setlength{\itemsep}{0pt}\setlength{\topsep}{2pt}}
\item[(A)] $1$
\item[(B)] $0$
\item[(C)] $3\pi$
\item[(D)] nonexistent
\end{list}

\bigskip\noindent\textbf{T22.}\ Let $f(x) = \begin{cases} \dfrac{x^2-4}{x-2} & \text{if } x \neq 2 \\ 7 & \text{if } x = 2 \end{cases}$

Which of the following statements is (are) true?

I. $\displaystyle\lim_{x\to 2} f(x)$ exists
II. $f(2)$ exists
III. $f$ is continuous at $x = 2$
\begin{list}{}{\setlength{\leftmargin}{3em}\setlength{\labelwidth}{2.5em}\setlength{\itemsep}{0pt}\setlength{\topsep}{2pt}}
\item[(A)] I only
\item[(B)] II only
\item[(C)] I and II only
\item[(D)] I, II, and III
\end{list}

\bigskip\noindent\textbf{T23.}\ If $\begin{cases} f(x) = \dfrac{x^2+3x}{6x} & \text{for } x \neq 0 \\ f(0) = k \end{cases}$

and if $f$ is continuous at $x = 0$, then $k =$
\begin{list}{}{\setlength{\leftmargin}{3em}\setlength{\labelwidth}{2.5em}\setlength{\itemsep}{0pt}\setlength{\topsep}{2pt}}
\item[(A)] $-\dfrac{1}{2}$
\item[(B)] $\dfrac{1}{6}$
\item[(C)] $\dfrac{1}{2}$
\item[(D)] $1$
\end{list}

\bigskip\noindent\textbf{T24.}\ Suppose $\begin{cases} f(x) = \dfrac{4x(x-1)}{x^2-4x+3} & \text{for } x \neq 1, 3 \\ f(1) = -2 \\ f(3) = 5 \end{cases}$

Then $f(x)$ is continuous
\begin{list}{}{\setlength{\leftmargin}{3em}\setlength{\labelwidth}{2.5em}\setlength{\itemsep}{0pt}\setlength{\topsep}{2pt}}
\item[(A)] except at $x = 1$
\item[(B)] except at $x = 3$
\item[(C)] except at $x = 1$ or $3$
\item[(D)] at each real number
\end{list}

\bigskip\noindent\textbf{T25.}\ The graph of $f(x) = \dfrac{9}{x^2-4}$ has
\begin{list}{}{\setlength{\leftmargin}{3em}\setlength{\labelwidth}{2.5em}\setlength{\itemsep}{0pt}\setlength{\topsep}{2pt}}
\item[(A)] one vertical asymptote, at $x = 2$
\item[(B)] the $y$-axis as its vertical asymptote
\item[(C)] the $x$-axis as its horizontal asymptote and $x = \pm 2$ as its vertical asymptotes
\item[(D)] two vertical asymptotes, at $x = \pm 2$, but no horizontal asymptote
\end{list}

\bigskip\noindent\textbf{T26.}\ The graph of $y = \dfrac{3x^2+3x+5}{6x^2-6x}$ has
\begin{list}{}{\setlength{\leftmargin}{3em}\setlength{\labelwidth}{2.5em}\setlength{\itemsep}{0pt}\setlength{\topsep}{2pt}}
\item[(A)] a horizontal asymptote at $y = \dfrac{1}{2}$ but no vertical asymptote
\item[(B)] no horizontal asymptote but two vertical asymptotes, at $x = 0$ and $x = 1$
\item[(C)] a horizontal asymptote at $y = \dfrac{1}{2}$ and two vertical asymptotes, at $x = 0$ and $x = 1$
\item[(D)] a horizontal asymptote at $y = \dfrac{1}{2}$ and two vertical asymptotes, at $x = \pm 1$
\end{list}

\bigskip\noindent\textbf{T27.}\ Let $f(x) = \begin{cases} \dfrac{2x^2}{x} & (x \neq 0) \\ 0 & (x = 0) \end{cases}$

Which of the following statements is (are) true?

I. $f(0)$ exists
II. $\displaystyle\lim_{x\to 0} f(x)$ exists
III. $f$ is continuous at $x = 0$
\begin{list}{}{\setlength{\leftmargin}{3em}\setlength{\labelwidth}{2.5em}\setlength{\itemsep}{0pt}\setlength{\topsep}{2pt}}
\item[(A)] I only
\item[(B)] II only
\item[(C)] I and II only
\item[(D)] I, II, and III
\end{list}

\bigskip\noindent\textbf{TB1.}\ If $[x]$ denotes the greatest integer not greater than $x$, then $\displaystyle\lim_{x\to 3/2} [x]$ is
\begin{list}{}{\setlength{\leftmargin}{3em}\setlength{\labelwidth}{2.5em}\setlength{\itemsep}{0pt}\setlength{\topsep}{2pt}}
\item[(A)] $\dfrac{3}{2}$
\item[(B)] $1$
\item[(C)] $0$
\item[(D)] nonexistent
\end{list}

\bigskip\noindent\textbf{TB2.}\ If $[x]$ denotes the greatest integer not greater than $x$, then $\displaystyle\lim_{x\to -4} [x]$ is
\begin{list}{}{\setlength{\leftmargin}{3em}\setlength{\labelwidth}{2.5em}\setlength{\itemsep}{0pt}\setlength{\topsep}{2pt}}
\item[(A)] $-5$
\item[(B)] $-4$
\item[(C)] $-3$
\item[(D)] nonexistent
\end{list}

\bigskip\noindent\textbf{TB3.}\ $\displaystyle\lim_{x\to\infty} \cos x$
\begin{list}{}{\setlength{\leftmargin}{3em}\setlength{\labelwidth}{2.5em}\setlength{\itemsep}{0pt}\setlength{\topsep}{2pt}}
\item[(A)] is $-1$
\item[(B)] is infinity
\item[(C)] is zero
\item[(D)] does not exist
\end{list}

\bigskip\noindent\textbf{TB4.}\ The function $f(x) = \begin{cases} \dfrac{3x^2}{x} & (x \neq 0) \\ 0 & (x = 0) \end{cases}$
\begin{list}{}{\setlength{\leftmargin}{3em}\setlength{\labelwidth}{2.5em}\setlength{\itemsep}{0pt}\setlength{\topsep}{2pt}}
\item[(A)] is continuous everywhere
\item[(B)] is continuous except at $x = 0$
\item[(C)] has a removable discontinuity at $x = 0$
\item[(D)] has an infinite discontinuity at $x = 0$
\end{list}

\bigskip\noindent\textbf{TB5.}\ \[
g(x) = \begin{cases}
2-x & (-1 \leq x < 0) \\
3x^2-3 & (0 \leq x \leq 1) \\
-x+5 & (1 < x < 2) \\
5 & (x = 2) \\
3x-3 & (2 < x \leq 3)
\end{cases}
\]
$\displaystyle\lim_{x\to 2} g(x)$
\begin{list}{}{\setlength{\leftmargin}{3em}\setlength{\labelwidth}{2.5em}\setlength{\itemsep}{0pt}\setlength{\topsep}{2pt}}
\item[(A)] equals $5$
\item[(B)] equals $4$
\item[(C)] equals $3$
\item[(D)] does not exist
\end{list}

\bigskip\noindent\textbf{TB6.}\ \[
g(x) = \begin{cases}
2-x & (-1 \leq x < 0) \\
3x^2-3 & (0 \leq x \leq 1) \\
-x+5 & (1 < x < 2) \\
5 & (x = 2) \\
3x-3 & (2 < x \leq 3)
\end{cases}
\]
The function $g$ is defined on $[-1,3]$
\begin{list}{}{\setlength{\leftmargin}{3em}\setlength{\labelwidth}{2.5em}\setlength{\itemsep}{0pt}\setlength{\topsep}{2pt}}
\item[(A)] if $x \neq 0$
\item[(B)] if $x \neq 1$
\item[(C)] if $x \neq 2$
\item[(D)] at each $x$ in $[-1,3]$
\end{list}

\bigskip\noindent\textbf{TB7.}\ \[
g(x) = \begin{cases}
2-x & (-1 \leq x < 0) \\
3x^2-3 & (0 \leq x \leq 1) \\
-x+5 & (1 < x < 2) \\
5 & (x = 2) \\
3x-3 & (2 < x \leq 3)
\end{cases}
\]
The function $g$ has a removable discontinuity at
\begin{list}{}{\setlength{\leftmargin}{3em}\setlength{\labelwidth}{2.5em}\setlength{\itemsep}{0pt}\setlength{\topsep}{2pt}}
\item[(A)] $x = 0$
\item[(B)] $x = 1$
\item[(C)] $x = 2$
\item[(D)] $x = 3$
\end{list}

\bigskip\noindent\textbf{TB8.}\ \[
g(x) = \begin{cases}
2-x & (-1 \leq x < 0) \\
3x^2-3 & (0 \leq x \leq 1) \\
-x+5 & (1 < x < 2) \\
5 & (x = 2) \\
3x-3 & (2 < x \leq 3)
\end{cases}
\]
On which of the following intervals is $g$ continuous?
\begin{list}{}{\setlength{\leftmargin}{3em}\setlength{\labelwidth}{2.5em}\setlength{\itemsep}{0pt}\setlength{\topsep}{2pt}}
\item[(A)] $-1 \leq x \leq 0$
\item[(B)] $0 < x < 1$
\item[(C)] $1 \leq x \leq 2$
\item[(D)] $2 \leq x \leq 3$
\end{list}

\bigskip\noindent\textbf{TB9.}\ \[
g(x) = \begin{cases}
2-x & (-1 \leq x < 0) \\
3x^2-3 & (0 \leq x \leq 1) \\
-x+5 & (1 < x < 2) \\
5 & (x = 2) \\
3x-3 & (2 < x \leq 3)
\end{cases}
\]
The function $g$ has a jump discontinuity at
\begin{list}{}{\setlength{\leftmargin}{3em}\setlength{\labelwidth}{2.5em}\setlength{\itemsep}{0pt}\setlength{\topsep}{2pt}}
\item[(A)] $x = -1$
\item[(B)] $x = 1$
\item[(C)] $x = 2$
\item[(D)] $x = 3$
\end{list}

\bigskip\noindent\textbf{TB10.}\ $\displaystyle\lim_{x\to 0} \sqrt{5+\arctan\dfrac{1}{x}}$
\begin{list}{}{\setlength{\leftmargin}{3em}\setlength{\labelwidth}{2.5em}\setlength{\itemsep}{0pt}\setlength{\topsep}{2pt}}
\item[(A)] $=\sqrt{5-\dfrac{\pi}{2}}$
\item[(B)] $=\sqrt{5+\dfrac{\pi}{2}}$
\item[(C)] $=\infty$
\item[(D)] does not exist
\end{list}

\bigskip\noindent\textbf{TB11.}\ Suppose $\displaystyle\lim_{x\to 5^-} f(x) = 4$, $\displaystyle\lim_{x\to 5^+} f(x) = 4$, and $f(5)$ is not defined. Which of the following statements is (are) true?

I. $\displaystyle\lim_{x\to 5} f(x) = 4$
II. $f$ is continuous everywhere except at $x = 5$
III. $f$ has a removable discontinuity at $x = 5$
\begin{list}{}{\setlength{\leftmargin}{3em}\setlength{\labelwidth}{2.5em}\setlength{\itemsep}{0pt}\setlength{\topsep}{2pt}}
\item[(A)] I only
\item[(B)] III only
\item[(C)] I and III only
\item[(D)] I, II, and III
\end{list}

\bigskip\noindent\textbf{TB12.}\ If $y = \dfrac{1}{4+10^{1/x}}$, then $\displaystyle\lim_{x\to 0} y$
\begin{list}{}{\setlength{\leftmargin}{3em}\setlength{\labelwidth}{2.5em}\setlength{\itemsep}{0pt}\setlength{\topsep}{2pt}}
\item[(A)] $=0$
\item[(B)] $=\dfrac{1}{14}$
\item[(C)] $=\dfrac{1}{4}$
\item[(D)] does not exist
\end{list}

\bigskip\noindent\textbf{TB13.}\ \[
h(x) = \begin{cases}
-x-2 & (-3 \leq x < -1) \\
2 & (x = -1) \\
x & (-1 < x < 1) \\
-x+4 & (1 \leq x < 2) \\
-x+5 & (2 \leq x \leq 3)
\end{cases}
\]
For what value(s) of $a$ is it true that $\displaystyle\lim_{x\to a} h(x)$ exists and $h(a)$ exists, but $\displaystyle\lim_{x\to a} h(x) \neq h(a)$? It is possible that $a =$
\begin{list}{}{\setlength{\leftmargin}{3em}\setlength{\labelwidth}{2.5em}\setlength{\itemsep}{0pt}\setlength{\topsep}{2pt}}
\item[(A)] $-1$ only
\item[(B)] $2$ only
\item[(C)] $-1$ or $1$ only
\item[(D)] $-1$ or $2$ only
\end{list}

\bigskip\noindent\textbf{TB14.}\ \[
h(x) = \begin{cases}
-x-2 & (-3 \leq x < -1) \\
2 & (x = -1) \\
x & (-1 < x < 1) \\
-x+4 & (1 \leq x < 2) \\
-x+5 & (2 \leq x \leq 3)
\end{cases}
\]
$\displaystyle\lim_{x\to a} h(x)$ does not exist for $a =$
\begin{list}{}{\setlength{\leftmargin}{3em}\setlength{\labelwidth}{2.5em}\setlength{\itemsep}{0pt}\setlength{\topsep}{2pt}}
\item[(A)] $-1$ only
\item[(B)] $1$ only
\item[(C)] $2$ only
\item[(D)] $1$ and $2$ only
\end{list}

\bigskip\noindent\textbf{TB15.}\ \[
h(x) = \begin{cases}
-x-2 & (-3 \leq x < -1) \\
2 & (x = -1) \\
x & (-1 < x < 1) \\
-x+4 & (1 \leq x < 2) \\
-x+5 & (2 \leq x \leq 3)
\end{cases}
\]
Which of the following statements about limits at $x = 2$ is (are) true?

I. $\displaystyle\lim_{x\to 2^-} h(x)$ exists
II. $\displaystyle\lim_{x\to 2^+} h(x)$ exists
III. $\displaystyle\lim_{x\to 2} h(x)$ exists
\begin{list}{}{\setlength{\leftmargin}{3em}\setlength{\labelwidth}{2.5em}\setlength{\itemsep}{0pt}\setlength{\topsep}{2pt}}
\item[(A)] I only
\item[(B)] II only
\item[(C)] I and II only
\item[(D)] I, II, and III
\end{list}

\end{document}
